% Post-CP24 isolated candidate. No triple-free theta image is claimed.
\subsection{An explicit quaternion norm-seven interpolation problem}
\status{Candidate; geometric fibre bound remains open.}

Let the ground field be algebraically closed of characteristic two.
For $r\ne0,1$, choose $t,\beta,w$ satisfying
\[
 z=\frac{r^2+1}{r^3},\quad D=r^2=t^2+t,\quad
 a=z^2,\quad b=z^2+z+1,\quad
 \beta^2+\beta=z+1,\quad w^2+w=at^5+bt^3.
\]
Write $C:y^2+y=F(x)=ax^5+bx^3$, and put
\[
 G(x)=zx^2+(z+1)x+\beta,\qquad
 H(x)=ux^2+vx+w,
 \quad u=z(t+r),\quad v=zt^2+t+r^{-2}.
\]
Then
\[
 g(x,y)=(x+1,y+G(x)),\qquad
 h(x,y)=(x+t,y+H(x))
\]
are automorphisms with $g^2=h^2=\iota$ and $gh=\iota hg$.
Indeed the three coefficient identities
\[
 u^2=at,\qquad v^2+u=bt,\qquad v=at^4+bt^2
\]
verify the curve equation, while
\[
 ut^2+vt=1,\qquad zt^2+(z+1)t+u+v=1
\]
verify the square and commutator. Their actions on the Jacobian
therefore form a quaternion group. Rosati sends $g,h,gh$ to their
negatives. Consequently
\[
 \psi=2+g+h+gh,\qquad \psi^\dagger\psi=7,
 \qquad\deg\psi=49
\]
is an etale polarized self-isogeny of multiplier seven.
Every nonzero $z$ admits the chosen $r$, by the cubic
$zr^3+r^2+1=0$; its roots are neither zero nor one.
The already proved $C_z$ presentation therefore covers every
two-rank-zero genus-two Jacobian. This gives an explicit norm-seven
self-isogeny, but says nothing yet about its theta fibres.
The existence of abstract noncommuting order-four partners is also
contained in \citep[Lemma~2.1 and proof, printed pp.~4--8]{BrockHowe2025Richelot}. That source's odd-norm nodality
is not asserted or imported here.

For $P=(p,y_0)$, set $X=x+p$, $d_1=F'(p)=ap^4+bp^2$, and
\[
 c_g=zD+t,\qquad c=Dc_g,
 \qquad Z=H(p)+t^2G(p)+Dd_1+tc_g,
 \qquad E=tG(p)+H(p)+(t+1)c_g.
\]
Define
\[
 A_3=\frac{ZE+cc_g}{D},\qquad
 A_2=(Z+c)G(p)+cd_1+A_3,\qquad A_1=cd_1,
\]
and the function
\[
 f_P=(ZX+c)(y+y_0)+A_1X+A_2X^2+A_3X^3.
\]
It vanishes along $2P+gP+hP+ghP$. Its norm under the
hyperelliptic involution factors exactly as
\begin{align*}
 f_P\,\iota(f_P)
 &=X^2(X+1)(X+t)(X+t+1)\,R_P(X),\\
 R_P(X)
 &=aZ^2X^2+(A_3^2+apZ^2)X
       +(b+a(D+1))Z^2+ac^2.
\end{align*}
Thus away from leading-coefficient degenerations this is a
degree-two residual effective divisor representing $-\psi(P-W)$.
All leading-coefficient and special-divisor cases must be handled
before using it to infer a fibre bound.

There is a useful linearized identity. Put
\[
 U_4(p)=p^4+(D+1)p^2+Dp,\qquad
 Z_0=w+t^2\beta+tc_g.
\]
Then
\[
 Z=aD\,U_4(p)+p^2+Z_0,\qquad
 Z(p+1)=Z(p)+1,\qquad Z(p+t)=Z(p)+t^2.
\]
The global constant coefficient of the residual polynomial, in
$x=X+p$, is
\[
 pA_3^2+(b+a(D+1))Z^2+ac^2.
\]

There is a further useful reduction, still submitted for audit.
Replacing the chosen root $t$ by $t+1$ changes $c_g=zD+t$ by one,
so choose $t$ with $c_g\ne0$. For any finite $P$, an arbitrary
section of $H^0(7W)$ vanishing on
$D_P=2P+gP+hP+ghP$ has the form
\[
 (LX+C)(y+y_0)+Cd_1X+B_2X^2+B_3X^3,
\]
where its three remaining value equations are equivalent to
\[
 CZ=LDc_g,\qquad
 B_3=(LE+Cc_g)/D,\qquad
 B_2=(L+C)G(p)+Cd_1+B_3.
\]
Because $Dc_g\ne0$, this solution space has dimension exactly
one for every finite $P$. Hence the residual effective degree-two
divisor $E_P\sim7W-D_P$ is nonspecial. In particular
$\psi(P-W)\ne0$: otherwise $D_P\sim5W$ and
$h^0(7W-D_P)=h^0(2W)=2$, a contradiction.
This proves $\ker\psi\cap\Theta=\{W\}$ and excludes tangent
identifications of a hyperelliptic pair.

The polynomial $Z$ has derivative $aD^2\ne0$ and degree four.
At one of its four roots, $A_3=c_g^2\ne0$, so $f_P$ has pole
order six and $E_P=W+R_P$. The displayed norm formula gives
\[
                  x(R_P)=p+\frac{aD^2}{c_g^2}.
\]
Different roots of $Z$ therefore give distinct residual
$x$-coordinates. The two points with a fixed $p$ give distinct
opposite images, since the image is nonzero and the Jacobian has
no nonzero geometric two-torsion. Thus these leading-coefficient
degenerations cannot supply a triple fibre. They also cannot
coincide with an image for $Z\ne0$, whose nonspecial reduced
divisor is finite of degree two.

For the remaining open set $Z\ne0$, put
\[
 E_0=w+t\beta+(t+1)c_g,
 \qquad q_0=(E_0^2+Z_0)/D^2.
\]
The additional exact identity
\[
 E^2+apD^2=Z+E_0^2+Z_0
\]
shows that the monic residual linear coefficient is
\[
 \ell=p+\frac{A_3^2}{aZ^2}
      =\frac{Z}{aD^2}+\frac{q_0}{a}
         +\frac{c_g^4}{aZ^2}.
\]
It is a separable degree-three rational function of $Z$.
This is the precise remaining issue: equality of the residual
quadratics alone permits a cubic ambiguity. The ordinate
coefficient must also be controlled to exclude three distinct
residual representations. No such exclusion is claimed here.

The retained program
\path{code/nonordinary_seven_q8_interpolation.py} verifies seventeen
identities over $\mathbf F_2(r)$ modulo the three displayed
Artin--Schreier equations. The first execution passed at
20:59:00.912911 UTC in $0.2739150630077347$ seconds, with full output
\path{outputs/nonordinary_seven_q8_interpolation.json}.
That exact seventeen-identity version and output are retained
separately in the history and output directories. A later version
adds the eighteenth identity expressing $\ell$ through $Z$.
No theta nodality or degree-seven monodromy conclusion follows
from this certificate alone. A subsequent rational-coordinate
exploration first failed because the symbolic coefficient domain
did not implement a field gcd; the exact version and traceback are
retained. Its second version was interrupted before a result and
does not supply a completed calculation.
